\documentclass[11pt]{article}
\usepackage[utf8]{inputenc}
\usepackage{amsmath, amssymb, amsthm}
\newtheorem{theorem}{Theorem}[section]

\title{\(1+\delta\) approximation in sublinear time}
\author{Ian Chi}
\date{}

\begin{document}

\title

\section{Definitions}

For a function \(f:[n]\to \mathbb R\) we say a pair \((i,i+1)\) for \(i\) in \([n-1]\) is 
\begin{enumerate}
    \item a \emph{good pair} if there exists an optimal LIS that does not change the values in this pair
    \item a \emph{bad pair} if no optimal LIS includes this pair untouched
    \item a \emph{perfect pair} is a pair that must remain untouched in all optimal LIS
\end{enumerate}

Note that all \((j,j+1)\) such that \(f(j)>f(j+1)\) are "bad pairs". Also note that all perfect pairs are also good pairs 


\section{\(1+\delta\) assuming \(\varepsilon_f<0.5\)}

\begin{theorem}
    For a fixed delta, if the true \(\varepsilon_f\) is less than \(0.5\), Using distance-separation\((f,\varepsilon,\delta)\), the decision problem of if \(f\) is within \(\varepsilon\pm \delta\) is solvable in polylog time. 
\end{theorem}

\section{good pair analysis}

\subsection{fraction of good pairs}

\begin{theorem}
    At least half of all increasing pairs are good pairs
\end{theorem}

\subsection{padding good pairs}

\begin{theorem}
    given a good pair \((i,i+1)\), constructing a new function \(f_p:[n+1]\to \mathbb R\) were \begin{equation}
        f_p(x) = \begin{cases}
            f(x),& x\leq i\\
            \frac{f(i)+f(i+1)}{2},& x=i+1\\
            f(x+1),&x\geq i+2
        \end{cases}
    \end{equation}
    forces \((i,i+1),(i+1,i+2)\) to both be perfect pairs in the new function \(f_p\)
    
    % adding an additional value between \(i,i+1\) that respects monotonicity results in \((i,i+1),(i+1,i+2)\) to become perfect pairs. 
\end{theorem}




\end{document}